\chapter{State of the Art}
\label{ch:grundlagen}


\section{Term Definition}

\begin{description}
    [font=\normalfont,leftmargin=1.9in,style=multiline]
    \item[SURF IA]
        Enhanced Traffic Situational Awareness on the Airport Surface with Indications and Alerts, used to describe a system that provides alerts to pilots regarding traffic on the runway. (SURF A - only alerts, SURF IA - both indications and alerts)
    \item[Ownship]
        Ownship is to be defined as the aircraft equipped with the hypothetical SURF IA system.
    \item[Target]
        Target is to be defined as the aircraft or ground vehicle that is considered traffic by ownship.
    \item[Conflict Geometry]
        Parameter set that results in an accident if no evasive action is taken.
    \item[Convergence]
        Decreasing distance between ownship and traffic toward a common location on the runway.
        
        
    \item[Runway Geometry]
        Runway configuration in relation to another runway or taxiway.
\end{description}


\subsection{Error Causes}

\begin{description}
    [font=\normalfont,leftmargin=1.9in,style=multiline]
        \item[Pilot Deviation]
        Also called Pilot Error, an action by the pilot that deviates from intended operations, aviation regulations and procedures, or from air traffic controller commands.
    \item[Operational Incident]
        Also called Operational Error, an action by air traffic controllers that result in less than the minimum required separation between aircraft and obstacles, or an aircraft trespassing on a runway delegated to another aircraft's operations.
    \item[Vehicle/Pedestrian Deviation]
        The unauthorized presence of a ground vehicle or a person on a runway delegated to another aircraft's operations that interferes with the aircraft's operations.
\end{description}

Pilot Deviation is the leading cause of runway incursions. Landing and takeoff are the phases of flight where pilot workload is at its maximum, and this high level of workload can often cause loss of situational awareness or a breakdown of communications. The pilot can also experience the urge to do something while being instructed to wait and holding points, leading to violations of aviation regulations and actions that put the pilot and aircraft at risk. Communication breakdowns are also partly responsible for pilot errors, where a pilot receives a landing clearance intended for another aircraft or a pilot fails to correctly readback the ATC command.

Operational Incidents are the second leading cause of runway incursion. Handling multiple planes and/or runways at a time can place a high level of stress on the controller, causing them to momentarily forget about an aircraft or clearance issued on an active runway. The controller can also misjudge the separation of aircraft on approach or ensure that the command is received. ICAO lists other factors that might affect the controller personally such as distraction, workload, experience level, and the human machine interface.

Vehicle/Pedestrian Deviation is the least most common type of runway incursions, and is often due to communication breakdown or lack of familiarization with the aerodrome.

\subsection{Detection Types}

\begin{description}
    [font=\normalfont,leftmargin=1.9in,style=multiline]
    \item[True Detection]
        Scenario when an alert or indication is needed and provided.
    \item[Untrue Detection]
        Case when an alert or indication is unneeded but provided, caused by algorithmic errors or sensor failures. Untrue detections can be false, caused by a failure of the alerting system, or nuisance, caused by algorithmic errors.
    \item[Missed Detection]
        Case when an alert or indication is needed but not provided, caused by algorithmic errors or sensor failures.
\end{description}

\section{Current Runway Incursion Mitigation Measures}

The aviation industry have already taken numerous measures to reduce the risk of runway incursions, and currently all of them are airport based. These range from infrastructure installed in the airport, such as runway safety lights and ground radar, to operational procedures on ensuring proper separation on taxiing, takeoff, and landing around runways.

\subsection{Lighting Solutions}

Runway lighting is an effective visual indicator to show if the runway is occupied or not, with a high visibility of up to 5 miles in daytime and approximately 20 miles at night-time.

Runway Status Lights (RWSL), Stopbars, Runway Guard Lights (RGL), and Final Approach Runway Occupancy Signal (FAROS) are lighting solutions designed to minimize the risk of runway incursions. RWSL works by using data from surveillance sources, such as airport surveillance radar, ground radar, and Multilateration (MLAT), generates an aerial or ground track, and predicts the most likely path and automatically determines the lighting configuration. Runway Entrance Lights (REL) and Runway Intersection Lights (RIL) turn red when a runway is unsafe to enter or cross, Takeoff Hold Lights (THLs) turn red when a runway is unsafe and causes the pilots to stop the aircraft, regardless of whether they have begun aircraft roll or not. A lit Stopbar requires the pilots to immediately stop and hold at the lit Stopbar and can only proceed when the Stopbar is switched off and the ATC gives the command to proceed. RGLs are flashing lights located where a taxiway connects to a runway and exist to increase a pilot's situational awareness, but a pilot does not need to take immediate action. Finally, FAROS automatically detects if an aircraft enters a runway and causes the landing guidance lights to go from steady to flashing.

\subsection{Operational Procedures}

Aside from runway lighting, air traffic control (ATC) also play an important part in runway incursion risk reduction. ATC controllers act off of available traffic information in order to provide guidance, commands, and advisories. Ensuring aircraft have proper separation is not only important in reducing the risk of runway incursions, but also ensures that aircraft landing or taking off do not fly into each other's wake turbulence. Wake turbulence can cause loss of speed and control, and presents a high hazard given the low altiltudes.

An action that minimizes communication breakdown and ensures correct reception of message is the read-back, a procedure where the receiver repeats the message or part of the message to obtain confirmation of correct reception. Some messages always require a read-back, such as ATC route clearances, any action involving any runway, level, heading, and speed instructions, and Special Service Request codes.

\subsection{Advanced Surface Movement Guidance and Control Systems}

Advanced Surface Movement Guidance and Control Systems, or A-MGCS, takes data from multiple sources, such as primary radar, MLAT sensors, surveillance radar, aircraft ADS-B transponders, and flight plan data, to provide enhanced surveillance, routing, guidance, and control functions, particularly in adverse meteorological or visibility conditions. Most well known implementation of A-MGCS is the FAA's ASDE-X, installed in 35 major airports in America, and it is primarily responsible for reducing the RIS rating of runway incursions at American airports.

ICAO defines 4 basic functions that an A-SMGC system should perform, the most basic of which is surveillance. The A-SMGCS system should provide positional data for all movement in the area of coverage, even in adverse meteorological conditions, and the identify and tag each movement. It should be able to provide routing, or designate a route for each aircraft or vehicle within the movement area and allow for destination or route modification. It should provide guidance to each pilot and ground vehicle driver for their designated path, as well as clearly indicate routes or areas which are unavailable or restricted. The final and highest function it should be able to perform is control, including provision of longitudinal separation based on type and state, providing hazard risk alerts, such as for runway incursions, and finally the ability to predict, detect, and resolve and conflict scenarios.



\section{SURF-IA}

SURF-IA aims to reduce the risk of runway incursions in the same way TCAS reduces the risk of airborne collisions, by providing the flight crew with indications and alerts and improving the flight crew's situational awareness. The system uses ADS-B transmit and receive functions to predict traffic paths and convergence rate based on ownship and traffic states. RTCA provides the minimum Operational Safety, Performance, and Interoperability Requirements (SPR) for SURF IA implementation in DO-323. 

\subsection{Equipment Requirements}

SURF-IA requires ADS-B at the very minimum, as it can only detect traffic equipped with ADS-B Out. Preferably, Traffic Information Service - Broadcast (TIS-B), which enables non-ADS-B equipped aircraft to be detected by other aircraft, and Automatic Dependent Surveillance - Rebroadcast (ADS-R), will allow ownship to detect all traffic. ADS-R is a ground station which receives all ADS-B data and rebroadcasts it to other aircraft in the vicinity.

\subsection{Data Requirements}

SURF-IA requires traffic states in order to calculate convergence rates, and the set of data requirements include an aircraft identifier and state data, including altitude, speed, and directionality and the accuracies thereof, as well as aircraft dimensions. For traffic targets, it requires the traffic category, which differentiates between aircraft, helicopters, and surface vehicles. Although helicopters and surface vehicles qualify as SURF-IA targets, they are not intended for use on either.

SURF-IA also requires airport map information, which includes the airport ID, runway parameters to determine runway boundaries, field elevation and accuracy of runway parameters. Other parameters such as locations of hold lines, runway shoulders, parallel taxiways are not required but can reduce the rate of nuisance alerts.

\subsection{Detection Logic}

SURF-IA area of coverage extends to 1,000 ft above field elevation and 3 NM away from the runway threshold. RTCA defines 6 aircraft states in DO-323: 1. Taxiing on Taxiway toward Hold-line or stopped at Hold-line, 2. Entering/Crossing/Exiting (runway), 3. Takeoff (from detection of takeoff roll to lift-off), 4. Approach (within area of coverage), 5.  After landing (speed greater or equal to 40 knots), and 6. Stopped or taxiing along runway. SURF IA takes ownship and traffic states and calculates a convergence rate to determine if separation will pose a safety hazard and provide the appropriate alerts.  

\subsection{Display of Alerts and Indications}


There are two types of alerts, both of which are required to have an aural component, which are presented when ownship and traffic are converging at a high speed. Caution alerts require immediate flight crew awareness and subsequent flight path modifying response, while a Warning alert requires both immediate flight crew awareness and flight path modifying response. The caution alerts provide sufficient time for pilots to evaluate situation and make a decision on how to proceed.

SURF IA provides two types of indications: Runway Status Indications, which indicates that the flight crew should verify runway status before proceeding, and a Traffic Indication, which indicates that no immediate collision hazard with ownship exists but one can develop over time.

[insert photo here]

SURF IA builds on two other RTCA systems for Cockpit Display of Traffic Information (CDTI): Airport Surface Situational Awareness (ASSA, DO-317), and Final Approach and Runway Occupancy Awareness (FAROA, DO-289). The alerts and indications appear on the flight displays, as well as ownship and relevant traffic in the vicinity. Traffic is considered relevant if it poses a collision hazard, or can develop into a collision hazard, and is highlight yellow or red depending on the level of hazard.

\subsection{SURF IA Implementations}

Currently, the only certified implementation of SURF IA is the ROA on Garmin's G1000 equipped Cessna Caravan. Garmin expects FAA certification for its G5000 avionics suite by late 2024, and although it is designed for transport category aircraft, Garmin is only currently targeting business jets.

\section{Review of Relevant Literature}

As one of ICAO's five global high risk categories, runway incursions have been the subject of much research.

\subsection{Research on Runway Incursions}

modeling runway incursion severity
- different types of conflict points
- airport surface geometry
\newline
runway incursion prevention: a technology solution \newline
- assumes inadvertent mistakes by pilot and/or controller and research hypotheses is that improvement of pilot and controller situational awareness reduces likelihood of mistakes \newline
- suggests tech solutions, not limited to ads b \newline
- notes that ground based uplink is not recommended as latency and error issues \newline
\newline
ripas (runway incursion prevention and avoidance system) \newline
- mature to level where reliable detection of ri \newline
- provides solution to non towered aircraft \newline
- highlights downside of surf ia - aircraft and vehicles non cooperative to surveillance \newline
- suggests update on infrastructure \newline
\newline 
RIS risk analysis - DOT 9717 DSI \newline
- OE incidents tend to be more severe \newline
- incident types and severity vary across regions \newline
- commercial aircraft tend to be involced in less severe accidents \newline
- controller workload has correlartion w/ probabilities of severe incidents \newline
\newline
risk based framework for assessment of ri events \newline
- severity vs risk based framework \newline
- severity - limited by uncontrolled random circumstances \newline
- risk based - inventory based \newline
- factors in this inventory include conflict geometries, aircraft types and states \newline
- proposes risk based assessment \newline




\subsection{Research on SURF IA}

Mostly focuses on human factors regarding hmi/cdti
- human factors flight test evaluation of surf ia
- flight simulation evaluation of surf ia
- using cdti to prevent RI
\newline
ERAU study
- very similar to this study
- shows that surf ia only showed 70 percent effective at triggering caution or warning alerts for icao serious model
- names 3 non alerting aircraft states which prevent caution or warning alert for icao serious model
- may be explained by risc is outcome based, takes into account closest proximity and maneuvers after closest proximity
- surf ia is collision risk based

